\chapter{The European Map}\label{ch:nw:the-european-map}

On 6 June 2022 Euronext finished moving its Core Data Centre, and the \omterm{def:nw:colocation-products-and-how-they-are-sold:colo}{colocation} that went with it, from Basildon, east of London, to
Aruba's IT3 data centre near Bergamo, in northern Italy. The building it left was ICE's; the building it moved to handles, in the words
of Euronext's chief executive, a quarter of European trading volumes. From Slough, where Euronext had put one of its London points
of presence, light in a vacuum had needed \qty{258}{\micro\second} to reach Basildon; it now needs
\qty{3.31}{\milli\second} to reach Bergamo. The firms that trade Euronext's markets from London started that week four and a half
milliseconds of straight glass further away, and two years later Euronext itself was selling them a microwave link to close part of the gap.

Europe has no single New Jersey. Its trading sites are spread over London, Frankfurt, Zurich, northern Italy, Stockholm and Paris, and the
routes between them (London to Frankfurt above all) are the continent's latency races. This chapter builds the European half of the site
table with the same rules as chapter~10, adds the registry that says which venue matches in which building on which date, and measures what
the migration did.

\section{The London-area sites}

\begin{definition}[Carrier-neutral data centre]\label{def:nw:the-european-map:neutral}
A \emph{carrier-neutral data centre}\index{carrier-neutral data centre} is one whose operator lets any telecom carrier bring its network into
the building and sells \omterm{def:nw:colocation-products-and-how-they-are-sold:mmr}{cross-connects} between any tenants and any carriers, rather than tying tenants to its own network or to one carrier.
\end{definition}

A venue can be in a carrier-neutral building or in its own; what matters to a firm is that it can reach the venue's systems, and its own
carriers, from its cabinet. Euronext describes IT3 as ``carrier-neutral by design'', with four carrier entry points (north, south, east and
west) and dark fibre to the Milan Internet eXchange; Deutsche Börse's co-location in Frankfurt is in a building where the housing contracts
are between the firm and the operator, Equinix, and the exchange is ``one of many customers''.

London has three sites for this chapter. Equinix's LD4, on the Slough Trading Estate west of London, is where Euronext put one of its two
London points of presence and where its microwave service to Bergamo starts; it is also one end of the London--Frankfurt routes below. The
London Stock Exchange moved its primary data centre from a site at Liverpool Street, in the City, to Telehouse's Docklands campus: planned
for October 2022, the cutover of Millennium Trading, Turquoise and the Group Ticker Plant was moved to 18 February 2023. ICE's European
Liquidity Centre, in Basildon, Essex, runs ICE Futures Europe, and until June 2022 it also hosted Euronext.

\begin{definition}[Point of presence]\label{def:nw:the-european-map:pop}
A \emph{point of presence}\index{point of presence} (PoP) of a network is a place, usually racks in a \omterm{def:nw:the-european-map:neutral}{carrier-neutral data centre}, where the
network's owner (a venue, a carrier) accepts customers' connections and carries their traffic to its own sites, so that a customer
can reach a distant venue with a \omterm{def:nw:colocation-products-and-how-they-are-sold:mmr}{cross-connect} in a building near it.
\end{definition}

The three London sites are close: 44~kilometres from Slough to Docklands, 34 from Docklands to Basildon, 77 from Slough to Basildon, that
is 147, 114 and \qty{258}{\micro\second} of light in vacuum. Inside such a triangle a firm can serve all three venues from
one building, with a few hundred microseconds of glass to each; the migration took one corner of it to Italy.

\begin{dated}{2026-09}{Who matches where in Europe}\label{dat:nw:the-european-map:sites}
\begin{itemize}
\item \textbf{Bergamo} (Aruba Global Cloud Data Centre IT3, Via San Clemente 53, Ponte San Pietro): Euronext's Core Data Centre since 6 June
  2022, with its \omterm{def:nw:colocation-products-and-how-they-are-sold:colo}{colocation}, run by Euronext; the disaster-recovery site stays in the Paris region; London points of presence at InterXion
  and Slough LD4.
\item \textbf{Frankfurt} (Equinix FR2, Kruppstrasse 121--127): the primary Eurex and Xetra back-ends and Deutsche Börse's co-location.
\item \textbf{Zurich} (Equinix ZH4, Josefstrasse 225): SIX Swiss Exchange's co-location, managed by SIX; a SIX microwave network links ZH4
  to FR2.
\item \textbf{London Docklands} (Telehouse campus, Coriander Avenue): the London Stock Exchange Group's primary data centre since the
  migration from the City.
\item \textbf{Basildon} (ICE's European Liquidity Centre): ICE Futures Europe; the street is from a directory and the site is approximate.
\item \textbf{Upplands Väsby} (Smedbyvägen 6), 25~kilometres from Stockholm: Nasdaq's Nordic primary data centre; its secondary moved to
  Equinix SK2, south of Stockholm, in 2025.
\end{itemize}
\end{dated}

\section{Frankfurt and Zurich}

Deutsche Börse's page is precise: the Equinix FR2 data centre in Frankfurt is the co-location data centre ``in which the primary Eurex and
Xetra back-ends are located''. Its co-location connections are 10~gigabits, with transaction and market data on separate connections so that
they do not queue behind each other (chapter~1), priced by the month: 6\,000~EUR for one Xetra market-data or order-entry connection, 8\,400
for two market-data products together, and 400~EUR for its PTP time service. Chapter~9's American fee schedules are filed with the SEC; the
European ones are published as price lists, and they are just as volatile.

SIX runs its co-location in Equinix ZH4 in Zurich and publishes an average latency of \qty{10.6}{\micro\second} for a co-located 10-gigabit
connection to its matcher. It also runs a microwave network from ZH4 to FR2, at 10~megabits a second, with equal access for all trading
participants: an exchange selling the corridor to a neighbouring market's building, as Euronext does from London. Frankfurt and Zurich are
\qty{306.7}{\kilo\metre} apart, \qty{1.02}{\milli\second} in vacuum and \qty{1.50}{\milli\second} in straight fibre.

\begin{table}[ht]
\centering\small
\begin{tabular}{llrrr}
\toprule
From & To & Geodesic & Vacuum & Fibre \\
 & & (km) & (\si{\micro\second}) & (\si{\micro\second}) \\
\midrule
Slough LD4 & Frankfurt FR2 & 677.7 & 2\,260.6 & 3\,305.0 \\
Slough LD4 & Bergamo IT3 & 992.2 & 3\,309.7 & 4\,838.7 \\
Frankfurt FR2 & Bergamo IT3 & 497.6 & 1\,659.9 & 2\,426.8 \\
Slough LD4 & Basildon & 77.4 & 258.0 & 377.2 \\
Docklands & Basildon & 34.0 & 113.5 & 166.0 \\
Slough LD4 & Docklands & 44.0 & 146.9 & 214.8 \\
Basildon & Frankfurt FR2 & 603.2 & 2\,012.2 & 2\,941.8 \\
Basildon & Bergamo IT3 & 936.2 & 3\,122.9 & 4\,565.7 \\
Frankfurt FR2 & Zurich ZH4 & 306.7 & 1\,023.1 & 1\,495.8 \\
Zurich ZH4 & Bergamo IT3 & 204.5 & 682.2 & 997.4 \\
Frankfurt FR2 & Upplands Väsby & 1\,195.0 & 3\,986.0 & 5\,827.6 \\
Slough LD4 & Upplands Väsby & 1\,462.5 & 4\,878.2 & 7\,132.0 \\
\bottomrule
\end{tabular}
\caption{Distances and one-way \omterm{def:nw:the-north-american-map:floor}{latency floors} between the European sites of \cref{dat:nw:the-european-map:sites} (WGS-84 geodesics; fibre
at $n_g = 1.4620$). Basildon, Bergamo and Upplands Väsby are placed at street level, a kilometre or two from the buildings, which moves their
floors by a few microseconds. Data: \texttt{nw\_europe.pair\_rows()}.}\label{tab:nw:the-european-map:pairs}
\end{table}

\begin{omfigure}
\begin{tikzpicture}
\begin{axis}[axis equal image,width=11.5cm,xmin=-560,xmax=1060,ymin=-440,ymax=380,xlabel={km east of 4.5\textdegree{} E},
  ylabel={km north of 49\textdegree{} N},tick label style={font=\footnotesize,/pgf/number format/1000 sep={\,}},label style={font=\small},grid=major,grid style={gray!20},
  table/col sep=comma,unbounded coords=jump]
\addplot[omDef,thick,no markers] table[x=x,y=y]{figdata/networks/11-the-european-map/eu_links.csv};
\addplot[only marks,mark=*,mark size=2.4pt,omThm,visualization depends on={value \thisrow{anchor}\as\anc},nodes near coords,
  point meta=explicit symbolic,every node near coord/.append style={font=\footnotesize,anchor=\anc,black}]
  table[x=x,y=y,meta=name]{figdata/networks/11-the-european-map/eu_sites.csv};
\addplot[only marks,mark=none,nodes near coords,point meta=explicit symbolic,
  every node near coord/.append style={font=\scriptsize,fill=white,inner sep=1pt,anchor=center,text=omMeth}]
  table[x=x,y=y,meta=text]{figdata/networks/11-the-european-map/eu_labels.csv};
\end{axis}
\end{tikzpicture}

\omcaption{London, Frankfurt, Zurich and Bergamo from computed coordinates, each link labelled with its \omterm{def:nw:the-north-american-map:geodesic}{geodesic distance} and light-time in
vacuum: the London--Frankfurt--Bergamo triangle that Euronext's migration created. A local equirectangular projection centred on
49\textdegree{}~N, 4.5\textdegree{}~E, for drawing only. Data: \texttt{fig\_europe.py} on \texttt{firm.geomap}.}\label{fig:nw:the-european-map:map}
\end{omfigure}

\section{The northern-Italy migration}

Euronext's frequently asked questions about the move, as it stood in November 2021, describe the old arrangement and the new one. Euronext
relied on ICE's Basildon data centre under a long-term service agreement; \omterm{def:nw:colocation-products-and-how-they-are-sold:colo}{colocation} was, in its press release's word, ``outsourced''.
It chose a third-party building in Bergamo, in the country of Borsa Italiana, which it had just bought, to use the experience of the
team that ran Borsa Italiana's \omterm{def:nw:colocation-products-and-how-they-are-sold:colo}{colocation}. It kept its disaster-recovery site in the Paris region. It opened two points of presence in London, at
InterXion and at Slough LD4, from which a client could reach Bergamo, the disaster-recovery site and Borsa Italiana's Milan data centre,
but not Basildon. And it opened the roof: wireless providers and trading firms could install antennas on a dedicated area of the roof, with
Euronext guaranteeing equalisation from the back of the antenna to the \omterm{def:nw:colocation-products-and-how-they-are-sold:mmr}{meet-me room}, the rule of chapter~9 applied to the air.

\begin{omfigure}
\begin{tikzpicture}[font=\footnotesize,b/.style={draw,rounded corners=2pt,minimum height=0.6cm,align=center,inner sep=3pt},>=Stealth]
\node[b,fill=omDef!15] (fc) at (-0.8,0) {firm's cabinet\\ in Slough LD4};
\node[b,fill=omThm!15] (pop) at (3.6,0) {Euronext\\ point of presence\\ (LD4)};
\node[b,fill=gray!10,minimum width=3.1cm] (net) at (7.6,0) {Euronext network\\ (carriers' fibre)};
\node[b,fill=omThm!25] (me) at (11.3,0.9) {Bergamo IT3:\\ matching engines};
\node[b,fill=omThm!10] (dr) at (11.3,-0.9) {Paris region:\\ disaster recovery};
\draw[->,thick] (fc) -- node[above,font=\scriptsize]{cross-connect} (pop);
\draw[->,thick] (pop) -- (net);
\draw[->,thick] (net.east) -- (me.west);
\draw[->,thick] (net.east) -- (dr.west);
\node[b,fill=omMeth!12] (mw) at (3.6,1.9) {microwave (EWIN)\\ antennas};
\draw[->,thick,omMeth,dashed] (fc.north) |- (mw.west);
\draw[->,thick,omMeth,dashed] (mw.east) -| node[pos=0.25,above,font=\scriptsize]{less than 4 ms, fibre back-up} (me.north);
\node[font=\scriptsize,align=center,gray] at (7.6,-1.6) {a customer in another building reaches the PoP by a carrier's circuit};
\end{tikzpicture}

\omcaption{How a London firm reaches Euronext after the migration, schematically: a \omterm{def:nw:colocation-products-and-how-they-are-sold:mmr}{cross-connect} in LD4 to Euronext's \omterm{def:nw:the-european-map:pop}{point of presence},
Euronext's network to Bergamo and to the disaster-recovery site, or, since 2024, Euronext's microwave service from LD4 with fibre back-up.
Sources: Euronext's migration FAQ and its 2024 release (\cref{dat:nw:the-european-map:sites}).}\label{fig:nw:the-european-map:pop}
\end{omfigure}

\begin{omfigure}
\begin{tikzpicture}
\begin{axis}[ybar,width=11.5cm,height=6.2cm,bar width=7pt,ymin=0,ymax=5.4,enlarge x limits=0.15,
  symbolic x coords={Slough LD4,Docklands,FR2,ZH4},xtick=data,ylabel={one-way floor to Euronext (ms)},
  tick label style={font=\footnotesize},label style={font=\small},grid=major,grid style={gray!25},legend columns=4,
  legend style={font=\footnotesize,at={(0.5,-0.18)},anchor=north,draw=none,/tikz/every even column/.append style={column sep=6pt}},table/col sep=comma]
\addplot[fill=omDef!30,draw=omDef] table[x=from,y=vac_before]{figdata/networks/11-the-european-map/migration.csv};
\addplot[fill=omDef!80,draw=omDef] table[x=from,y=vac_after]{figdata/networks/11-the-european-map/migration.csv};
\addplot[fill=omThm!25,draw=omThm] table[x=from,y=fib_before]{figdata/networks/11-the-european-map/migration.csv};
\addplot[fill=omThm!75,draw=omThm] table[x=from,y=fib_after]{figdata/networks/11-the-european-map/migration.csv};
\legend{vacuum before,vacuum after,fibre before,fibre after}
\end{axis}
\end{tikzpicture}

\omcaption{The one-way \omterm{def:nw:the-north-american-map:floor}{latency floor} from four sites to Euronext's primary data centre before (Basildon, to 5 June 2022) and after (Bergamo,
from 6 June), in vacuum and in straight fibre. London moved three milliseconds away in vacuum, four and a half in fibre; Frankfurt came
slightly closer and Zurich much closer. Data: \texttt{nw\_europe.migration()} on \texttt{firm.venuesites}.}\label{fig:nw:the-european-map:move}
\end{omfigure}

\Cref{fig:nw:the-european-map:move} is the migration in numbers. From Slough the floor in vacuum went from \qty{258}{\micro\second} to
\qty{3310}{\micro\second}, \qty{3.05}{\milli\second} more; in fibre from \qty{377}{\micro\second} to \qty{4839}{\micro\second},
\qty{4.46}{\milli\second} more. From Docklands, much the same. From Frankfurt the floor fell by \qty{352}{\micro\second} (2.01 to
\qty{1.66}{\milli\second}), and from Zurich by \qty{1.81}{\milli\second}. The London firm that had traded Euronext's markets against
London's, a few hundred microseconds apart, now trades them against each other from opposite ends of a three-millisecond corridor, and
chooses where to put the servers that react to each.

\omcode{code/networks/11-the-european-map/python/nw_europe.py}{53}{62}{The migration's floors: the registry's primary site for Euronext Paris
on the day before and the day of the move, and the floors to each from a list of sites.}\label{lst:nw:the-european-map:move}

\begin{proposition}[Moving the engine moves the triangle]\label{prop:nw:the-european-map:triangle}
Let a firm trade two venues at sites $A$ and $B$ from a third site $F$. The time a signal takes from one venue to the other through the
firm is at least $d(A,F) + d(F,B)$ divided by the speed of the medium, and at least $d(A,B)$ over it along the direct route; the firm's
position matters only through the first sum, which is smallest when $F$ lies on the geodesic between $A$ and $B$.
\end{proposition}

\begin{proof}
The triangle inequality for geodesics: $d(A,F) + d(F,B) \ge d(A,B)$, with equality exactly when $F$ lies on the shortest path between $A$ and
$B$.
\end{proof}

Before the migration, London, Frankfurt and Euronext formed a sliver: 678, 603 and 77~kilometres, with Euronext's engine almost on top of
London. After it, the triangle is wide: 678, 498 and 992~kilometres. A firm that relays between Frankfurt and Bergamo through London pays
$678 + 992 = 1\,670$ kilometres of path for a direct distance of 498; a firm that wants to be between Euronext and Eurex should sit on the
Frankfurt--Bergamo line, which is where Zurich almost lies (307 plus 205~kilometres, 511 against 498).

\section{The Nordics and the rest}

Nasdaq's Nordic primary data centre stands in Upplands Väsby, 25~kilometres north of Stockholm, and its secondary moved in 2025 to Equinix
SK2, south of the city. It is \qty{1462.5}{\kilo\metre} from Slough (\qty{4.88}{\milli\second} in vacuum) and \qty{1195}{\kilo\metre} from
Frankfurt: the farthest of the chapter's sites from London, and a site to which chapters~22 and~23 return.

The rest of Europe follows the same pattern and the same method: each venue's documents name a building, often in a \omterm{def:nw:the-european-map:neutral}{carrier-neutral data
centre} in its own city; the registry holds one row per venue, role and period, with its source. Euronext's 2024 release speaks of a single
liquidity pool on one trading platform covering seven European markets, with its core data centre in Bergamo; the registry records there
the six markets that the migration's sources name, and leaves Borsa Italiana's until a source dates its move.

\omcode{code/firm/venuesites/firm_venuesites.py}{78}{98}{The venue-to-site registry's queries: rows by venue, role and date; a venue's primary
site on a date (an error if the table gives none or several); the venues matching in a building; and every venue's site by distance from
a given building.}\label{lst:nw:the-european-map:registry}

\begin{omfigure}
\begin{tikzpicture}
\begin{axis}[xbar,width=11cm,height=6cm,bar width=6pt,xmin=0,xmax=7.8,ytick=data,
  yticklabels from table={figdata/networks/11-the-european-map/from_ld4.csv}{site},y dir=reverse,
  xlabel={one-way latency floor from Slough LD4 (ms)},tick label style={font=\footnotesize},label style={font=\small},
  grid=major,grid style={gray!25},legend pos=north east,legend style={font=\footnotesize},table/col sep=comma]
\addplot[fill=omDef!55,draw=omDef] table[y expr=\coordindex,x=vacuum_ms]{figdata/networks/11-the-european-map/from_ld4.csv};
\addplot[fill=omThm!40,draw=omThm] table[y expr=\coordindex,x=fibre_ms]{figdata/networks/11-the-european-map/from_ld4.csv};
\legend{vacuum,fibre}
\end{axis}
\end{tikzpicture}

\omcaption{Every European venue site of the registry by its one-way floor from Slough LD4, in vacuum and in straight fibre: the London sites
within a third of a millisecond, Frankfurt and Zurich at two to four, Bergamo and Upplands Väsby beyond. Data: \texttt{fig\_europe.py},
\texttt{Registry.nearest} from LD4.}\label{fig:nw:the-european-map:ld4}
\end{omfigure}

\section{The London--Frankfurt corridor}

\begin{definition}[Latency corridor]\label{def:nw:the-european-map:corridor}
A \emph{latency corridor}\index{latency corridor} is a pair of trading sites between which several networks compete to carry
latency-sensitive traffic, each selling its route by its one-way latency: London--Frankfurt, Chicago--New Jersey, Zurich--Frankfurt.
\end{definition}

The corridor's floor is set by the geodesic: \qty{677.7}{\kilo\metre} from Slough LD4 to Frankfurt FR2, \qty{2.261}{\milli\second} in air and
\qty{3.305}{\milli\second} in straight fibre. In March 2015 McKay Brothers announced a route between the two buildings at
\qty{4.64}{\milli\second} round trip, \qty{2.32}{\milli\second} one way: a \omterm{def:nw:the-north-american-map:factor}{route factor} of 1.026 in air, \qty{59}{\micro\second} above the
floor, over a path that must cross the English Channel. It was already a millisecond faster than straight glass could ever be.

\begin{table}[ht]
\centering\small
\begin{tabular}{lrrrrr}
\toprule
Microwave route & Geodesic & Published & Floor in & Route & Fibre floor \\
 & (km) & (ms) & air (ms) & factor & (ms) \\
\midrule
LD4 -- FR2 (2015) & 677.7 & 2.32 & 2.261 & 1.026 & 3.305 \\
LD4 -- Bergamo IT3 (2024) & 992.2 & $< 4$ & 3.311 & $< 1.21$ & 4.839 \\
\bottomrule
\end{tabular}
\caption{Published European microwave routes against their floors: McKay Brothers' 2015 round trip halved, and Euronext's 2024 figure for
EWIN, which is an upper bound, so its \omterm{def:nw:the-north-american-map:factor}{route factor} is too. Data: \texttt{nw\_europe.route\_rows()}.}\label{tab:nw:the-european-map:routes}
\end{table}

The London--Bergamo corridor is Euronext's own. Its EWIN service, launched in July 2024 with McKay Brothers, carries orders by microwave from
LD4 to IT3 in ``less than 4 milliseconds'', with full fibre back-up. The floor in air is \qty{3.311}{\milli\second}, so the published bound
allows a \omterm{def:nw:the-north-american-map:factor}{route factor} of up to 1.21; straight fibre could not do better than \qty{4.839}{\milli\second}. A published bound is not a
measurement: it tells a firm what the venue will stand behind, and the firm measures the rest (chapter~15).

\begin{remark}[Orders one way, data the other]\label{rem:nw:the-european-map:direction}
EWIN carries orders from London to Bergamo; the firm still needs Euronext's market data in London, by fibre from the points of presence or by
any other route it buys. A corridor is two races, one per direction, and a firm may buy them from different networks.
\end{remark}

\begin{method}[Adding a venue to the registry]\label{met:nw:the-european-map:add}
\begin{enumerate}
\item Find the venue's MIC in the ISO 10383 list, and its site in the venue's own connectivity documents or filings.
\item Add the building to the site table with its address source and geocode (chapter~10's method), and a precision note if the geocode is
  only street-level.
\item Add one registry row per role (primary, secondary, disaster recovery, \omterm{def:nw:the-european-map:pop}{point of presence}, market-data distribution) with its source
  and the date read; a migration is two rows, the old one closed the day before the new one opens.
\item Recompute distances and floors, and date the whole table: venues move, as Euronext and the London Stock Exchange did within a year.
\end{enumerate}
\end{method}

\omcode{code/firm/venuesites/firm_venuesites.py}{51}{55}{A registry row's validity on a date: an open start says nothing about earlier dates;
an open end means the row still holds.}\label{lst:nw:the-european-map:dates}

\section{Tutorial: the European registry}\label{tut:nw:the-european-map:tut}

\begin{tutorial}
\textbf{Goal.} Extend the site table to Europe, build the venue-to-site registry, and measure Euronext's migration.
\textbf{End state:} \cref{fig:nw:the-european-map:map,fig:nw:the-european-map:move,fig:nw:the-european-map:ld4} and
\cref{tab:nw:the-european-map:pairs,tab:nw:the-european-map:routes}.
\begin{tutsteps}
\item \textbf{Sites.} Seven European rows in \texttt{data/networks/sites.csv}, each with its ledger rows; three are street-level and say so.
\item \textbf{Registry.} \texttt{code/firm/venuesites/data/venue\_sites.csv}: one row per venue, role and period;
  \texttt{Registry.load()} refuses a row without a source, with an unknown site, or with its dates reversed.
\item \textbf{The migration.} \texttt{nw\_europe.migration()} (\cref{lst:nw:the-european-map:move}) asks the registry for Euronext's primary
  site on 5 and on 6 June 2022 and computes the floors; \texttt{Registry.moves} lists the move itself for XPAR.
\item \textbf{Joining Book~1.} \texttt{join\_venues} pairs the registry with One Quant Book~1's \texttt{firm.venues} sample by MIC, without
  editing it; MICs absent from the sample are skipped, not invented.
\item \textbf{Figures.} \texttt{python fig\_europe.py} writes the map, the migration bars and the distances from LD4.
\end{tutsteps}
\textbf{What to change next.} Add Borsa Italiana's Milan data centre and the Paris-region disaster-recovery site once their buildings are
sourced; ask \texttt{Registry.nearest} from Frankfurt instead of Slough.
\end{tutorial}

\section{Build: the venue-to-site registry}\label{bld:nw:the-european-map:venuesites}

\begin{build}
\bfield{Purpose} Which venue matches, keeps its disaster recovery, and accepts connections in which building, on which date, with sources:
the input of the connectivity chapters~22 to~26 and of the plan of chapter~29.
\bfield{Interface} \texttt{firm\_venuesites}: \texttt{VenueSite(mic, venue, role, site, valid\_from, valid\_to, source, as\_of)},
\texttt{Registry.load}, \texttt{rows(mic, role, on)}, \texttt{primary(mic, on)}, \texttt{venues\_at(site, on)},
\texttt{nearest(site, role, on, medium)}, \texttt{moves(mic)}, \texttt{stale(today, days)}, \texttt{join\_venues(registry, venues)}.
\bfield{Rules} Every row has a source and a read date; every site is in \texttt{firm.geomap}'s table; a venue has exactly one primary site on
any date or the query fails; Book~1's registry is read, never edited.
\bfield{Acceptance tests} \texttt{code/firm/venuesites/tests/}: the Euronext migration by date, the nearest venues from LD4 in order, the
join with Book~1's sample, and each malformed row refused.
\bfield{Stretch} Read the site table and the registry from one versioned file with a history of changes; add each venue's market-data
distribution points and their \omterm{def:nw:networking-for-trading:group}{multicast groups}.
\end{build}

\begin{omsources}
\item Euronext, releases of 15 June 2022 (Core Data Centre migration) and 11 July 2024 (EWIN); \emph{Frequently Asked Questions about the
Euronext Data Centre Migration}, v6.1 (November 2021).
\item Deutsche Börse, Xetra co-location services; SIX, co-location service fees (January 2026) and connectivity page; Nasdaq, Primary
Datacenter Stockholm (Väsby) and the 2025 secondary-site notice; ICE, European Liquidity Centre operating policies; LSEG, Data Centre Migration
plan (2022).
\item McKay Brothers, release of 10 March 2015; ISO 10383 market identifier codes; site coordinates from OpenStreetMap Nominatim.
\end{omsources}

\section{Exercises}

\begin{exercise}[$\star$]\label{exo:nw:the-european-map:1}
By how much did the one-way floor in fibre from Slough LD4 to Euronext's primary data centre change on 6 June 2022?
\end{exercise}

\begin{exercise}[$\star$]\label{exo:nw:the-european-map:2}
McKay Brothers' 2015 figure for LD4--FR2 was a round trip of \qty{4.64}{\milli\second}. What is its \omterm{def:nw:the-north-american-map:factor}{route factor} in air?
\end{exercise}

\begin{exercise}[$\star$]\label{exo:nw:the-european-map:3}
What is the difference between a \omterm{def:nw:the-european-map:pop}{point of presence} and \omterm{def:nw:colocation-products-and-how-they-are-sold:colo}{colocation}?
\end{exercise}

\begin{exercise}[$\star\star$]\label{exo:nw:the-european-map:4}
Why did Euronext's FAQ promise equalisation ``from the back of the antenna to the Meet Me Room''? What would happen without it?
\end{exercise}

\begin{exercise}[$\star\star$]\label{exo:nw:the-european-map:5}
A firm relays signals between Eurex and Euronext. Compare the path through Slough with the path through Zurich, and with the direct one.
\end{exercise}

\begin{exercise}[$\star\star$]\label{exo:nw:the-european-map:6}
EWIN is published as ``less than \qty{4}{\milli\second}''. What does that tell you, and what does it not?
\end{exercise}

\begin{exercise}[$\star\star\star$]\label{exo:nw:the-european-map:7}
\emph{Coding.} Add a row to the registry that gives Euronext Paris two primary sites on the same day, and ask for its primary site. What happens,
and why is that the right behaviour?
\end{exercise}

\begin{exercise}[$\star\star\star$]\label{exo:nw:the-european-map:8}
\emph{Find the flaw.} ``Since the migration, London firms are three milliseconds behind Italian firms on Euronext, so London market makers
cannot compete there.''
\end{exercise}

\section{Problem: Basildon to Bergamo}

\begin{problem}[{Weekend problem --- what Euronext's migration did to the map}]\label{pb:nw:the-european-map:1}
A market maker quotes Euronext's markets from Slough LD4 and hedges on Eurex in Frankfurt FR2. On 6 June 2022 Euronext's Core Data Centre moves
from Basildon to Bergamo.

\medskip\noindent\textbf{Part I --- Before.}
\begin{enumerate}
\item What were the \omterm{def:nw:the-north-american-map:geodesic}{geodesic distance} and the one-way floors (vacuum, fibre) from LD4 to Basildon?
\item From Basildon to FR2?
\item What was the shape of the LD4--FR2--Euronext triangle?
\item Where could the firm put one server to be near both Euronext and London's other venues?
\end{enumerate}

\medskip\noindent\textbf{Part II --- After.}
\begin{enumerate}[resume]
\item What are the distances and floors from LD4 to Bergamo, in vacuum and in fibre?
\item By how much did each floor change?
\item From FR2 to Bergamo, and by how much did that change?
\item What is the triangle's new shape?
\end{enumerate}

\medskip\noindent\textbf{Part III --- The firm's choices.}
\begin{enumerate}[resume]
\item The firm keeps its servers in LD4 and reaches Bergamo over fibre with a \omterm{def:nw:the-north-american-map:factor}{route factor} of 1.2. What is its one-way latency?
\item It buys EWIN instead. At most how long do its orders take, and what does it save against part III's first answer?
\item Should its quoting engine for Euronext's markets move to Bergamo? What does it lose if it does?
\item How far from the Frankfurt--Bergamo geodesic is Zurich, as a detour?
\end{enumerate}

\medskip\noindent\textbf{Part IV --- The verdict.}
\begin{enumerate}[resume]
\item State the \emph{named result}: the change in the one-way floor from LD4 to Euronext's matching engines, and the new triangle.
\item Why did Euronext open points of presence in London before the move?
\item What does the registry record, and why two rows rather than an edited one?
\item What other European move of 2022--2023 should the registry record?
\item What does \cref{prop:nw:the-european-map:triangle} say about where to put a relay between Eurex and Euronext?
\item Which of the chapter's sites are only street-level, and does it matter for the named result?
\item What would a hollow-core fibre (chapter~13) change on the London--Bergamo corridor?
\item In one sentence: what did the migration change for a London firm?
\end{enumerate}
\end{problem}

\section{Interview questions}

\begin{interviewq}[$\star$ \iqroles{developer, trader}]\label{iq:nw:the-european-map:1}
Where are the main European matching engines, and what is the London--Frankfurt \omterm{def:nw:the-north-american-map:floor}{latency floor}?
\end{interviewq}

\begin{interviewq}[$\star$ \iqroles{developer}]\label{iq:nw:the-european-map:2}
What is a \omterm{def:nw:the-european-map:pop}{point of presence}, and why would an exchange open one in another country?
\end{interviewq}

\begin{interviewq}[$\star\star$ \iqroles{developer}]\label{iq:nw:the-european-map:3}
How would you keep a table of which venue is in which data centre correct over time?
\end{interviewq}

\begin{interviewq}[$\star\star$ \iqroles{trader, developer}]\label{iq:nw:the-european-map:4}
An exchange announces it will move its matching engine to another country. What do you do, and in what order?
\end{interviewq}

\begin{interviewq}[$\star\star$ \iqroles{trader}]\label{iq:nw:the-european-map:5}
Why do exchanges sell microwave links to their own data centres?
\end{interviewq}

\begin{interviewq}[$\star\star\star$ \iqroles{developer, researcher}]\label{iq:nw:the-european-map:6}
Where would you place servers to trade Eurex against Euronext, and how would you check the choice?
\end{interviewq}
